\begin{abstract}
	En el contexto de Análisis de Datos y de Machine Learning,
	donde el manejo de datos de alta dimensión es cada vez más
	frecuente, la reducción de dimensionalidad es una técnica clave
	para mejorar la eficiencia computacional y, por ende, reducir los
	dañinos efectos de la conocida ``maldición de la dimensionalidad''.
	El teorema de \emph{Johnson-Lindenstrauss} es un resultado
	fundamental, que establece que un conjunto de datos en un espacio
	de alta dimensión puede ser proyectado en un espacio de menor
	dimensión con una distorsión mínima de las distancias euclídeas.
	Estas proyecciones permiten preservar relaciones geométricas
	esenciales, justificando el uso de técnicas de Análisis de Datos en
	un espacio de dimensión significativamente menor.
	Este trabajo final de curso explora las peculiaridades que el
	tratamiento de datos de alta dimensión introduce, junto con las
	desigualdades de concentración útiles para conseguir probar el
	Teorema de Johnson-Lindenstrauss en su versión probabilística.
	También se analizarán, mediante simulaciones, estas peculiaridades
	y las implicaciones prácticas de este resultado.
\end{abstract}
